% Section 3: the carrier floor.  Printed in full, including the three
% constructions the phase-1.5 manuscript discharged by citing an evidence
% pack (the rounding, the vertical-direction choice, the negative curl).
\section{The carrier floor}\label{sec:floor}

The selector in Definition~\ref{def:X1} makes the contributing carriers
geometrically special, and that specialness bounds the degrees available to
their HOMFLY rows from below. The bound is the engine of Section~\ref{sec:s7}:
every vanishing proved there is an instance of it.

\subsection{What the selector forces}

\begin{remark}[attribution of the constructions in this section]\label{rem:floorsources}
The phase-1.5 manuscript proved Theorem~\ref{thm:carrierfloor}(C) but discharged
three of its steps by citing an evidence pack of the counterpart bench: the
rounding of the polygon corners, the choice of a vertical direction with its
tangency count, and the exact negative curl. Those three are printed here in
full --- Lemma~\ref{lem:rounding}, Theorem~\ref{thm:carrierfloor}(B) and Lemma~\ref{lem:curl} ---
and no conclusion here depends on that pack. The direct floor count in
Theorem~\ref{thm:carrierfloor}(B) replaces its former degree-theory last step;
attribution is otherwise unchanged: the rational curl model and its affine
insertion, the dual-vector argument of
Lemma~\ref{lem:uniformrot}(ii), and the computation identifying the forbidden
crossing as positive are the counterpart's; Lemma~\ref{lem:uniformrot}(i), the mirror calculation in the proof of
Theorem~\ref{thm:carrierfloor}(C), and that theorem's assembly are this author's. Printing a proof is not
claiming it; it is taking responsibility for it.
\end{remark}

\subsection{Rounding}

\begin{lemma}[rounding]\label{lem:rounding}
Here $\rot$ is as in Definition~\ref{def:rot}.
Let $L$ be a closed polygon with corners $q_1,\dots,q_c$, and let $D$ be an
oriented diagram whose underlying plane curve is $L$ --- $L$ together with an
over/under assignment at each of its double points. Assume the principal turns
of $L$ all exist and are \emph{nonzero}; that $L$ has finitely many double
points, all transversal, none of them a corner; and that no corner of $L$ lies
on an edge of $L$ other than the two incident to it. Then there is a \emph{clearance} $\varepsilon_0(L)>0$ such that for every
$\varepsilon\in\bigl(0,\varepsilon_0(L)\bigr)$ there is a $C^\infty$ regular
closed plane curve $L_\varepsilon$, and a diagram $D_\varepsilon$ carried by it, with
these properties. The clearance is part of the conclusion and is quantified
here: an earlier revision wrote "below the bound displayed in the proof", which
names an object chosen inside a proof as though the statement supplied it.

The hypothesis names a curve \emph{and} a diagram: clause~(c) speaks of an
over/under assignment, of crossing signs and of a writhe, none of which a plane
curve carries, and an earlier revision bound only $L$. It does \emph{not} name a
generic parent polygon. An earlier revision did, and then
Theorem~\ref{thm:carrierfloor}(C) applied this lemma on its own hypotheses, which
do not mention one --- the peer's polygon $((0,0),(1,0),(-2,-2),(2,0),(-2,1))$
satisfies the floor's hypotheses and is not a carrier of a support of a generic
polygon. What genericity was supplying is written out instead: finitely many
transverse double points, none at a corner, and no corner on a non-incident
edge, which is exactly what makes the finitely many distances in the proof
positive.
\begin{enumerate}
\item[(a)] $L_\varepsilon$ coincides with $L$ outside the union of the discs of
radius $\varepsilon$ about the corners.
\item[(b)] Inside the disc about $q_i$ the unit tangent moves \emph{strictly}
monotonically, in the sense of $\sgn\tau_i$, from $\delta_{i-1}/|\delta_{i-1}|$
to $\delta_i/|\delta_i|$, sweeping an arc of length exactly $|\tau_i|$ and no
more, and attaining each direction of that arc at exactly one parameter.
Moreover the unit-tangent map is an \emph{immersion} on the open junction arc:
parametrized by arclength, its angular derivative is nonzero at every interior
parameter, while at the two ends it and all its derivatives vanish, the
junction meeting the straight edges flat to infinite order.

Both clauses are exported because both are read downstream, and the second does
not follow from the first: a strictly monotone angle may still have a vanishing
derivative at a parameter, and the direct count below needs the derivative
itself to be nonzero, not merely the monotonicity. "Moves monotonically" alone
is weaker still --- it is
satisfied by a junction that pauses on one direction for a whole sub-arc, and
then a count of the parameters carrying a prescribed direction is infinite. Both
hold because the profile fixed in the proof below is strictly increasing on
$(0,1)$ with $\varphi'>0$ there, and flat to infinite order at $0$ and $1$: the
tangent angle at arclength $\sigma$ along the junction of length $\ell$ is
$\theta_u+\tau_i\varphi(\sigma/\ell)$, whose derivative is
$\tau_i\varphi'(\sigma/\ell)/\ell\neq0$ for $0<\sigma<\ell$.
\item[(c)] $L_\varepsilon$ has the same double points as $L$, with the same
strands, the same over/under assignment and hence the same crossing signs and
the same writhe.
\item[(d)] $\rot(L_\varepsilon)=\rot(L)$.
\item[(e)] \emph{The disc package is returned.} There are pairwise disjoint
closed discs $D_1,\dots,D_c$, one about each corner $q_i$, such that each $D_i$
meets no edge of $L$ that is not incident to $q_i$, contains no double point of
$L$, meets the two incident edges exactly in the two sub-segments of length
$\varepsilon$ at $q_i$, and contains the whole of the modification made at
$q_i$. This clause is a conclusion of the lemma, not a package handed over by a
paragraph of its proof; its consumer asks for a disc meeting the diagram in one
embedded arc and reads it here.
\end{enumerate}
\end{lemma}
\begin{proof}
\emph{The junction, constructed.} Since $0<|\tau_i|<\pi$, the unit directions
$u=\delta_{i-1}/|\delta_{i-1}|$ and $v=\delta_i/|\delta_i|$ are neither parallel
nor antiparallel. Write $A_0=q_i-\varepsilon u$ and $A_1=q_i+\varepsilon v$ for
the two points at distance $\varepsilon$ from $q_i$ along the incident edges, so
that the displacement to be realized is
\[
   A_1-A_0=\varepsilon\,(u+v),
\]
a nonzero vector along the bisector of the convex angle at $q_i$. Fix once and
for all the \emph{transition profile}
\[
   \varphi(t)=\frac{f(t)}{f(t)+f(1-t)},
   \qquad
   f(t)=\begin{cases} e^{-1/t}, & t>0,\\ 0,& t\leq0,\end{cases}
\]
which is exhibited rather than posited, and which has the four properties every
use below reads. It is $C^\infty$ on $[0,1]$: $f$ is $C^\infty$ on $\mathbb R$
with all derivatives vanishing at $0$, and the denominator
$f(t)+f(1-t)$ is positive on $[0,1]$, being $f(1)>0$ at $t=0$ and $f(1)>0$ at
$t=1$ and a sum of two nonnegative terms not both zero in between. It has
$\varphi(0)=0$ and $\varphi(1)=1$. It satisfies $\varphi(1-t)=1-\varphi(t)$, by
inspection of the formula. And it is \emph{strictly increasing on $(0,1)$}:
\[
   \varphi'(t)=\frac{f'(t)f(1-t)+f(t)f'(1-t)}{\bigl(f(t)+f(1-t)\bigr)^2}>0
   \qquad (0<t<1),
\]
both terms of the numerator being positive there, since $f>0$ and
$f'(t)=t^{-2}e^{-1/t}>0$ on $(0,1)$. Every derivative of $\varphi$ vanishes at
$0$ and at $1$, because every derivative of $f$ vanishes at $0$ and the
denominator is smooth and nonvanishing.

Strict increase is not a convenience. An earlier revision asked only that
$\varphi$ be nondecreasing, and a nondecreasing profile may be constant on a
subinterval; the tangent direction would then be constant on a whole sub-arc of
the junction, so a direction attained there would be attained at a continuum of
parameters rather than at isolated ones, and every count of the points carrying
a prescribed tangent direction --- the counts this section takes later --- would
be infinite. With $\varphi$ strictly increasing on $(0,1)$ the tangent angle
$\theta_u+\tau_i\varphi(t/\ell)$ is strictly monotone across the junction, and
each direction in the swept arc is attained exactly once. Let
$\theta_u$ be the angle of $u$ and, for a length $\ell>0$, let
\[
   \gamma_\ell(s)=A_0+\int_0^{s}
   \bigl(\cos(\theta_u+\tau_i\varphi(t/\ell)),\,
         \sin(\theta_u+\tau_i\varphi(t/\ell))\bigr)\,\mathrm dt,
   \qquad s\in[0,\ell].
\]
Then $\gamma_\ell$ is a unit-speed $C^\infty$ arc. Its tangent equals $u$
\emph{at} $s=0$ and $v$ \emph{at} $s=\ell$, and agrees with those constants
\emph{to infinite order} at those two parameters, every derivative of $\varphi$
vanishing there; it does not equal them on a neighbourhood, and cannot, since
the angle is strictly monotone on $(0,\ell)$ --- an earlier revision wrote
"near $s=0$", which contradicts the strictness this same construction exports.
Infinite-order agreement at the endpoint is what the junction needs: it is
exactly the condition under which replacing the corner by the arc leaves a
$C^\infty$ regular curve, since all one-sided derivatives match those of the
straight edge. The tangent angle moves strictly monotonically through exactly
$\tau_i$ and no more, and its derivative
$\tau_i\varphi'(s/\ell)/\ell$ is nonzero on $(0,\ell)$, which is (b).

Its endpoint is $A_0+\ell\,m(\varphi)$ where
$m(\varphi)=\int_0^1(\cos(\theta_u+\tau_i\varphi),\sin(\theta_u+\tau_i\varphi))\,
\mathrm dt$. The symmetry $\varphi(1-t)=1-\varphi(t)$ makes the angles
$\theta_u+\tau_i\varphi(t)$ and $\theta_u+\tau_i\varphi(1-t)$ reflections of
each other in the bisector, so $m(\varphi)$ points along $u+v$; and
$|m(\varphi)|>0$ because all the directions integrated lie in a closed angular
interval of width $|\tau_i|<\pi$. Taking
$\ell=\varepsilon\,|u+v|/|m(\varphi)|$ makes the endpoint exactly $A_1$. This is
where the circular arc of an earlier revision has been replaced: that arc is
$C^1$ and not $C^\infty$ at the two junctions, and the document then carried a
lemma to repair the loss.

\emph{The arc lies in the triangle $A_0q_iA_1$, in coordinates.} Put
$\alpha=|\tau_i|/2\in(0,\pi/2)$ and $s=\sgn\tau_i$, let $w$ be the unit vector
along $u+v$ and $z$ the unit vector with $\det(w,z)=1$, and take $q_i$ as
origin, writing a point as $(x_w,x_z)$ in that frame. Then
\[
   u=\cos\alpha\,w-s\sin\alpha\,z,
   \qquad
   v=\cos\alpha\,w+s\sin\alpha\,z,
\]
because $u$ and $v$ are unit vectors whose sum is along $w$ and whose angle is
$|\tau_i|=2\alpha$, turning from $u$ to $v$ by $\tau_i$. Hence
$A_0=-\varepsilon u$ has coordinates $(-\varepsilon\cos\alpha,\ s\varepsilon
\sin\alpha)$ and $A_1=\varepsilon v$ has $(\varepsilon\cos\alpha,\
s\varepsilon\sin\alpha)$, and the closed triangle with vertices $A_0,q_i,A_1$
is exactly
\begin{equation}
   T_i=\bigl\{\,(x_w,x_z):\ |x_w|\leq\cot\alpha\cdot(s\,x_z),
   \quad 0\leq s\,x_z\leq\varepsilon\sin\alpha\,\bigr\},
\label{eq:triangle}
\end{equation}
the first inequality being the pair of cone conditions at the apex $q_i$ and the
second the chord through $A_0$ and $A_1$: a point $\lambda(-u)+\mu v$ with
$\lambda,\mu\geq0$ has $x_w=(\mu-\lambda)\cos\alpha$ and
$s\,x_z=(\lambda+\mu)\sin\alpha$, so $|x_w|\leq(\lambda+\mu)\cos\alpha=\cot\alpha
\cdot(s\,x_z)$ with equality exactly on the two edges, and $\lambda+\mu\leq
\varepsilon$ is the chord.

The junction arc satisfies all three inequalities of \eqref{eq:triangle}. Write
its unit tangent as $T(\sigma)=\cos\beta(\sigma)\,w+s\sin\beta(\sigma)\,z$ with
$\beta(\sigma)=-\alpha+2\alpha\,\varphi(\sigma/\ell)$, which is the tangent
angle constructed above read in this frame: at $\sigma=0$ it is $u$
($\beta=-\alpha$), at $\sigma=\ell$ it is $v$ ($\beta=\alpha$), and $\beta$ is
strictly increasing with values in $[-\alpha,\alpha]$.
\emph{The two cone inequalities.} Let $n_1$ be the unit normal to $u$ with
$\langle n_1,v\rangle>0$. The line $\mathbb Ru$ contains both $q_i$ and $A_0$,
and $T_i$ lies in $\{\langle\cdot,n_1\rangle\geq0\}$. Along the arc,
$\tfrac{d}{d\sigma}\langle\gamma,n_1\rangle=\langle T,n_1\rangle\geq0$, because
$T=au+bv$ with $a,b\geq0$ --- a direction at angle $\beta\in[-\alpha,\alpha]$
from $w$ lies in the closed cone of $u$ and $v$ --- so
$\langle T,n_1\rangle=b\langle v,n_1\rangle\geq0$; and
$\langle\gamma(0),n_1\rangle=\langle-\varepsilon u,n_1\rangle=0$. Hence
$\langle\gamma,n_1\rangle\geq0$ throughout. Symmetrically, with $n_2$ the unit
normal to $v$ with $\langle n_2,A_0\rangle>0$, one has $\langle
T,n_2\rangle=a\langle u,n_2\rangle\leq0$, so $\langle\gamma,n_2\rangle$ is
nonincreasing and equals $\langle\varepsilon v,n_2\rangle=0$ at $\sigma=\ell$;
hence it is $\geq0$ throughout. Those two are the first inequality of
\eqref{eq:triangle}.
\emph{The chord inequality.} $s\,x_z(\sigma)=\varepsilon\sin\alpha+\int_0^\sigma
\sin\beta(t)\,dt$, and $\int_0^\sigma\sin\beta\leq0$ for every
$\sigma\in[0,\ell]$: $\beta$ is negative on $(0,\ell/2)$ and positive on
$(\ell/2,\ell)$, and the symmetry $\varphi(1-t)=1-\varphi(t)$ makes
$\beta(\ell-t)=-\beta(t)$, so the integral decreases to its minimum at
$\ell/2$ and increases back to $0$ at $\ell$. Hence $s\,x_z\leq\varepsilon
\sin\alpha$, which is the second. The lower bound $s\,x_z\geq0$ is implied by
the two cone inequalities, the apex being the origin.

So the arc lies in $T_i$, and $T_i$ lies in the closed disc $D_i$ of radius
$\varepsilon$ about $q_i$: the disc is convex and contains the three vertices,
$|A_0-q_i|=|A_1-q_i|=\varepsilon$. An earlier revision argued this by saying the
arc ``cannot leave the region cut off by the two tangent lines'', which names
the conclusion rather than proving it and leaves the chord side unaddressed.

For (a) and (c) put
\[
   \varepsilon_0(L)=\tfrac13\min\Bigl\{
 \min_{i}\min_{k\neq i}|q_k-q_i|,\;
 \min_{i}\operatorname{dist}\bigl(q_i,\,\textstyle\bigcup\{\text{closed edges of }L\text{ not incident to }q_i\}\bigr),\;
 \min_i|\delta_i|,\;
 \min_{i,\,x\text{ a double point}}|x-q_i|
\Bigr\},
\]
and require $\varepsilon<\varepsilon_0(L)$; the index $i$ is bound by the outer
minimum in every term, an earlier revision having left it free in the first.
Each minimum over an empty index set is $+\infty$ --- which is how the last
one is read when $L$ has no double point. Every listed quantity is positive: the
corners are finitely many and distinct; a corner is at positive distance from a
non-incident closed edge because the edges are compact and a corner on one would
contradict the hypothesis that no corner lies on a non-incident edge; the edges have positive length; and a double point lies
in the relative interiors of two edges, so it is distinct from every corner.
With this $\varepsilon$ the discs $D_i$ are pairwise disjoint, each $D_i$ meets
no edge not incident to $q_i$, each $D_i$ contains no double point, and each
$D_i$ meets the two incident edges in the two sub-segments of length
$\varepsilon$ at $q_i$, and each contains the whole modification made at its
corner, the replacement arc lying within distance $\varepsilon$ of $q_i$. That
is clause~(e). Hence the rounding changes the curve only inside those
discs, where it meets nothing else, so no double point is created or destroyed
and the strands through each surviving double point are unchanged as arcs, with
their orientations; the over/under assignment is inherited, so the crossing signs
and the writhe are unchanged. That is (a) and (c). For (d), the total turning of
$L_\varepsilon$ is the sum of the turnings of its arcs, which is
$\sum_i\tau_i$ because the tangent is constant along the straight parts; and the
rotation number of a $C^1$ regular closed curve is $\frac1{2\pi}$ times its
total turning (Lemma~\ref{lem:turnlift}(iii)), which is also the value assigned to $L$
by Definition~\ref{def:rot}.
\end{proof}

\begin{lemma}[rotation bounds for uniform and one-dissent polygons]
\label{lem:uniformrot}
Here $\rot$ is as in Definition~\ref{def:rot}.
\begin{enumerate}
\item[(i)] Let $L$ be a closed polygon all of whose principal turns exist and are
positive. Then $\rot(L)\geq1$, with equality if $L$ has three corners. If all
turns are negative, $\rot(L)\leq-1$, again with equality in absolute value for
three corners.
\item[(ii)] Let $L$ be a closed polygon whose principal turns all exist, with
exactly one negative principal turn, of magnitude $\alpha\in(0,\pi)$, and let
$\Pi$ be the sum of the positive ones. Then
$\Pi-\alpha=2\pi r$ with $r=\rot(L)$ an integer, and $r\geq1$.
\end{enumerate}
\end{lemma}
\begin{proof}
(i) By Lemma~\ref{lem:turnlift}(ii),
$2\pi\rot(L)=\sum_i\tau_i$. A polygon has at least three corners and each
$\tau_i>0$, so the sum is positive and $\rot(L)>0$; being an integer,
$\rot(L)\geq1$. With three corners the sum is also below $3\pi$, forcing
$\rot(L)=1$. Orientation reversal in the same lemma gives both negative
claims.

(ii) The identity is Lemma~\ref{lem:turnlift}(ii). Suppose $r\leq0$. Then
$\Pi\leq\alpha<\pi$. Cut the traversal of $L$ immediately after the negative
corner and follow it once around: the direction of the edge is turned by the
successive principal turns, all of which are \emph{nonnegative} --- the
hypothesis names one negative turn and admits turns equal to zero, which $\Pi$
omits, and an earlier revision said "all of which are positive", which is false
as written when a zero turn is present --- and they sum to $\Pi<\pi$, so every
edge direction of $L$ lies in a closed angular interval $I$ of width $\Pi<\pi$.
A zero turn does not widen that interval, so the argument is unchanged. An
interval of width less than $\pi$ has a nonempty open dual: there is
$u\in\mathbb R^2$ with $\langle u,\delta\rangle>0$ for every direction
$\delta\in I$, hence for every edge of $L$. But the edges of a closed polygon
sum to zero, so
$0=\langle u,\sum_i\delta_i\rangle=\sum_i\langle u,\delta_i\rangle>0$, a
contradiction. Hence $r\geq1$, and $R:=|\rot(L)|=r$.
\end{proof}

\subsection{The negative curl}

\begin{lemma}[exact negative-curl replacement]\label{lem:curl}
Here $\rot$ is as in Definition~\ref{def:rot}.
Let $F$ be a connected $C^\infty$ immersed circle in the plane --- one
component, with finitely many transverse double points and no triple points ---
given with an oriented diagram, and let $p$ be a point of $F$ at which the
tangent points in a fixed direction $u$, isolated among such points, lying in
an embedded arc of $F$ that contains no double point and along which the tangent
turns strictly positively. Then $F$ may be modified inside a disc $\Delta$
meeting the rest of the diagram only in that arc, so that the resulting
diagram $F'$
\begin{enumerate}
\item[(i)] is again such an oriented diagram, satisfies
$P_{F'}(a,z)=P_F(a,z)$, and has the same double points outside $\Delta$, with
the same signs;
\item[(ii)] has no point of $\Delta$ at which the tangent equals $u$, and exactly
one at which it equals $-u$;
\item[(iii)] has exactly one double point inside $\Delta$, and it is negative;
\item[(iv)] satisfies $\rot(F')=\rot(F)-1$ and $w(F')=w(F)-1$.
\end{enumerate}
\end{lemma}
\begin{proof}
Work first in the rational model on $-2\leq t\leq2$,
\[
b(t)=\Bigl(\tfrac{3(t^2+7)}{11},\,-3t\Bigr),\qquad
c(t)=\bigl(t^2-1,\;t-t^3\bigr).
\]
The endpoints and endpoint tangent rays agree up to positive scale:
$b(\pm2)=c(\pm2)=(3,\mp6)$ and $b'(\pm2)=\frac3{11}c'(\pm2)$. The old arc turns
strictly positively and has exactly one tangent pointing straight down, at
$t=0$:
\[
\det\bigl(b'(t),b''(t)\bigr)=\tfrac{18}{11}>0,\qquad b'(0)=(0,-3).
\]
The replacement turns strictly negatively and has exactly one tangent pointing
straight up, at $t=0$:
\[
\det\bigl(c'(t),c''(t)\bigr)=-2(1+3t^2)<0,\qquad c'(0)=(0,1).
\]
In each arc the first derivative has vanishing first component only at $t=0$, so
neither arc has any other vertical tangent; this is (ii) in the model. The
replacement runs between the same endpoint rays along the complementary,
clockwise tangent path, so its compatible lift increment is the old one minus
$2\pi$; Lemma~\ref{lem:turnlift}(iii) gives (iv) for $\rot$. It has exactly one double point,
\[
c(-1)=c(1)=(0,0),\qquad c'(-1)=(-2,-2),\qquad c'(1)=(2,-2),
\]
and putting the later branch over the earlier one makes it negative, since a
crossing sign is the sign of the determinant of the overpassing tangent followed
by the underpassing tangent and
$\det\bigl(c'(1),c'(-1)\bigr)=-8<0$; that is (iii), and with it (iv) for $w$.
Being a single Reidemeister-I monogon, the replacement does not change the knot
type, which is (i).

To insert the model at $p$ we first record the projective fact it needs, in the
form the insertion argument requires.

\emph{Ordered-ray lemma.} Let $(r_1,r_2,r_3)$ and $(s_1,s_2,s_3)$ be triples of
rays from the origin such that $r_2$ lies in the open convex cone spanned by
$r_1,r_3$ and $s_2$ lies in the open convex cone spanned by $s_1,s_3$, with
$r_1,r_3$ and $s_1,s_3$ independent, and suppose in addition that for positive
representatives
\begin{equation}
   \sgn\det(\rho_1,\rho_3)=\sgn\det(\sigma_1,\sigma_3),
\label{eq:orderedray}
\end{equation}
a condition independent of the representatives chosen. Then there is a linear
map $A$ with $\det A>0$ carrying each $r_k$ into $s_k$.
\emph{Proof:} if both determinants in \eqref{eq:orderedray} are negative,
exchange the names $1$ and $3$ \emph{in both triples simultaneously}; this
preserves the correspondence $r_k\mapsto s_k$ and makes both determinants
positive. The linear map $A_0$ with $A_0\rho_1=\sigma_1$, $A_0\rho_3=\sigma_3$
then has positive determinant and carries $r_1,r_3$ to $s_1,s_3$. It carries the
open cone of $r_1,r_3$ onto that of $s_1,s_3$, so $A_0r_2$ is a ray $\sigma$ in
the latter cone. Writing $s_2=\mathbb R_{>0}(x\sigma_1+y\sigma_3)$ and
$\sigma=\mathbb R_{>0}(x'\sigma_1+y'\sigma_3)$ with $x,y,x',y'>0$, the map
$D=\mathrm{diag}(x/x',\,y/y')$ in the basis $(\sigma_1,\sigma_3)$ has positive
determinant, fixes $s_1$ and $s_3$, and carries $\sigma$ to $s_2$. Take
$A=DA_0$. $\square$

An earlier revision of this lemma omitted \eqref{eq:orderedray} and said the
positive determinants could be arranged ``after swapping the names once if
necessary''. A swap performed in one triple only destroys the correspondence it
is supposed to produce, and the peer refuted that step; the hypothesis is now
stated and the swap is simultaneous. In the application below the condition is
not inferred from the two cones: both determinants are computed and are
positive.

\emph{A positive-turn chart, and the cuts.} Parametrise the embedded arc by
arclength as $\gamma(s)$ with $\gamma(0)=p$ and unit tangent $\gamma'(0)=u$.
Because the tangent turns strictly positively, after shrinking the chart there
is a continuous strictly increasing lift $\theta$ with
\[
   \gamma'(s)=\cos\theta(s)\,u+\sin\theta(s)\,Ju,\qquad
   \theta(0)=0,\qquad |\theta(s)|<\tfrac\pi2,
\]
where $J$ is the positive quarter turn. Put $v=-Ju$, so that $(v,u)$ is a
positively oriented orthonormal basis, and set
\[
   \xi(s)=\langle\gamma(s)-p,\,v\rangle,\qquad
   \eta(s)=\langle\gamma(s)-p,\,u\rangle,
\]
so that $\xi'(s)=-\sin\theta(s)$ and $\eta'(s)=\cos\theta(s)$. Hence $\xi$
increases strictly for $s<0$ and decreases strictly for $s>0$, with a strict
maximum $\xi(0)=0$, while $\eta$ increases throughout. For every level
sufficiently close to $\xi(0)$ from below, strict monotonicity and the
intermediate value theorem give unique cuts $s_-<0<s_+$ with
$\xi(s_-)=\xi(s_+)$, and both tend to $0$ as the level tends to $\xi(0)$. Write
$p_\pm=\gamma(s_\pm)$. Equal transverse coordinates and monotone $\eta$ give
\[
   p_+-p_-=\ell u,\qquad \ell>0,\qquad \ell\to0 .
\]
Two side facts follow at once and are what the earlier revision lacked: the
central arc $\gamma((s_-,s_+))$ lies strictly in $\xi>\xi(s_-)$, and both
retained tails lie strictly in $\xi<\xi(s_-)$.

\emph{The affine fit, quantitatively.} In the positively oriented model basis
$v_0=(-1,0)$, $u_0=(0,-1)$ one has $q=\tfrac4{11}$ and
\[
   b'(-2)\in\mathbb R_{>0}(qv_0+u_0),\qquad
   b'(0)\in\mathbb R_{>0}u_0,\qquad
   b'(2)\in\mathbb R_{>0}(-qv_0+u_0).
\]
Write the two target endpoint tangents in the basis $(v,u)$ as
$T_-=\gamma'(s_-)=av+bu$ and $T_+=\gamma'(s_+)=-cv+du$; then
$a=-\sin\theta(s_-)>0$, $c=\sin\theta(s_+)>0$ and $b,d=\cos\theta(s_\mp)>0$,
and $a,c\to0$, $b,d\to1$ as the cuts approach $p$. The two outer determinants
are $\det(qv_0+u_0,-qv_0+u_0)=2q>0$ and $\det(T_-,T_+)=ad+bc>0$, so
\eqref{eq:orderedray} holds; it is computed, not inferred. Define
\[
   H=\frac{2}{\tfrac ba+\tfrac dc}=\frac{2ac}{bc+ad},\qquad
   x=\frac Hq,\qquad
   y=\frac{\tfrac ba-\tfrac dc}{q\bigl(\tfrac ba+\tfrac dc\bigr)}
    =\frac{bc-ad}{q(bc+ad)},
\]
and let $A$ be the linear map with $A u_0=u$ and $Av_0=xv+yu$. Then
\[
   A(qv_0+u_0)=\tfrac Ha\,T_-,\qquad
   A(-qv_0+u_0)=\tfrac Hc\,T_+,\qquad
   \det A=x=\frac{2ac}{q(bc+ad)}>0,
\]
each by substitution, and
\[
   1-(qy)^2=\frac{4abcd}{(bc+ad)^2}>0,
\]
so $|y|<1/q$ uniformly, while $x\to0$ as the cuts approach $p$: the map $A$
stays bounded. This is the quantitative statement that the earlier revision
replaced by an unquantified ``fixed unit-normalized part of $A$''.

Since $b(2)-b(-2)=12u_0$, put $B=\tfrac{\ell}{12}A$ and
$\Phi(z)=p_-+B\bigl(z-b(-2)\bigr)$. Then $\Phi(b(-2))=p_-$,
$\Phi(b(2))=p_-+\ell A u_0=p_-+\ell u=p_+$, and the endpoint tangent rays match
with positive scale. Because $\ell\to0$ while $A$ stays bounded, the diameter of
$\Phi(c([-2,2]))$ tends to $0$, so the cuts may be chosen so that the whole
inserted arc lies in a preassigned disc $\Delta$ about $p$ whose intersection
with the diagram is contained in the embedded chart --- the disc is fixed first
and the cuts afterwards.

\emph{Only the model's own double point is created.} In the basis $(v_0,u_0)$
the transverse coordinate of $c(t)$ is $1-t^2$, which equals $-3$ at both
endpoints, so the model's transverse excess over its endpoint chord is
$4-t^2>0$ for $-2<t<2$. Since $Bv_0=\tfrac{\ell}{12}(xv+yu)$ with $x>0$ and
$Bu_0=\tfrac{\ell}{12}u$, the physical transverse excess of the inserted arc
over the chord $[p_-,p_+]$ is exactly
\[
   \frac{\ell x}{12}\bigl(4-t^2\bigr)>0 .
\]
So the inserted interior lies strictly in $\xi>\xi(s_-)$ while both retained
tails lie strictly in $\xi<\xi(s_-)$: they meet only at the two joins. Together
with the shrinking of the previous paragraph, which keeps the inserted arc away
from every nonlocal part of the diagram, this \emph{proves} rather than assumes
that the only double point created is the model's own.

\emph{The replacement, and why the curve stays $C^\infty$ regular.} Replace the
sub-arc of $F$ between $p_-$ and $p_+$ by $\Phi\circ c$, modified near its two
endpoints as follows. At each join the two one-sided unit tangents already
agree, $\Phi c'(\pm2)$ being a positive multiple of $\Phi b'(\pm2)$ and that of
the old tangent, so the concatenation is $C^1$; it need not be $C^\infty$, and
the collar that repairs that is constructed here rather than described. An
earlier revision said only that one may ``interpolate the tangent angle by the
same device'', which does not say what curve results, and interpolating angles
moves the endpoint.

\emph{The collar, as a graph.} Work at the join $p_-$ first; the join at
$p_+$ is constructed after property (3) below, with its own signs printed
rather than left to "roles exchanged". In the chart coordinates
$(\xi,\eta)$ of the positive-turn chart --- $\xi=\langle\gamma-p,v\rangle$ with
$v=-Ju$, $\eta=\langle\gamma-p,u\rangle$ --- the old arc is a graph
$\xi=f(\eta)$, since $\eta'=\langle\gamma',u\rangle=\cos\theta>0$ there. Its
slope carries a sign that an earlier revision got wrong, and the sign is
load-bearing below, so it is computed:
$\xi'=\langle\gamma',v\rangle=-\sin\theta$, whence
\[
   f'(\eta)=\frac{\xi'}{\eta'}=-\tan\theta .
\]
For $\eta<\eta(p)$ the angle $\theta$ is strictly negative, $\theta$ being
strictly increasing with $\theta(0)=0$, so $f'>0$ there --- not $f'<0$, which is
what the earlier revision's $f'=\tan\theta$ gave. Let
$\eta_-=\eta(p_-)<\eta(p)$. The inserted arc leaves $p_-$ with the same unit
tangent, so near its start it is also a graph $\xi=g(\eta)$ with
$g(\eta_-)=f(\eta_-)$ and $g'(\eta_-)=f'(\eta_-)$; and because the model turns
strictly negatively its tangent angle only decreases from $\theta(s_-)<0$, so by
the same formula
\[
   g'>f'>0
\]
on that stretch. Two consequences are used below: $g>f$ on
$(\eta_-,\eta_-+\lambda]$, the two agreeing at $\eta_-$ with $g'>f'$; and at
every point of the collar the smaller slope is $f'(\eta)$, still positive. The
bound is pointwise of necessity: $f'=-\tan\theta$ \emph{decreases} along the
collar, $\theta$ increasing, so $0<f'(\eta)<f'(\eta_-)$ there and an endpoint
bound from the left end is unavailable --- an earlier revision claimed both
slopes bounded below by $f'(\eta_-)$, which $f'$ itself violates. Fix
$\lambda>0$ small enough that $[\eta_-,\eta_-+\lambda]$ lies both below
$\eta(p)$ and inside the stretch on which the inserted arc is a graph, and set,
on $[\eta_-,\eta_-+\lambda]$,
\[
   h(\eta)=\bigl(1-\varphi_\lambda(\eta)\bigr)f(\eta)+\varphi_\lambda(\eta)\,
   g(\eta),
   \qquad
   \varphi_\lambda(\eta)=\varphi\Bigl(\frac{\eta-\eta_-}{\lambda}\Bigr),
\]
with $\varphi$ the transition profile of Lemma~\ref{lem:rounding}. The modified
curve runs along the old arc up to $\eta_-$, along the graph of $h$ across the
collar, and along the inserted arc from $\eta_-+\lambda$ on. Three properties,
each proved from the formula.

\emph{(1) It is $C^\infty$ and regular, with positive speed.} $h$ is a $C^\infty$
function, and every derivative of $\varphi$ vanishes at $0$ and $1$, so $h$
agrees with $f$ to infinite order at $\eta_-$ and with $g$ to infinite order at
$\eta_-+\lambda$: the two joins are $C^\infty$. Parametrized by $\eta$, the
collar has velocity $h'(\eta)\,v+u$ in the frame $(v,u)$, of norm
$\sqrt{1+h'(\eta)^2}\geq1$, so its speed is positive and it is regular; being a
graph, it is embedded.

\emph{(2) No tangent of the collar equals $u$, and no smallness condition is
needed.} The tangent is parallel to $u$ exactly where $h'=0$, and
\[
   h'=(1-\varphi_\lambda)f'+\varphi_\lambda g'
      +\frac{\varphi'\bigl((\eta-\eta_-)/\lambda\bigr)}{\lambda}\,
       \bigl(g-f\bigr).
\]
Every term is nonnegative, and the sum of the first two is positive: it is the
convex combination $(1-\varphi_\lambda)f'+\varphi_\lambda g'$ of two positive
numbers --- each \emph{term} separately can vanish, $\varphi_\lambda$ being $0$
at the left end and $1$ at the right, which an earlier revision's "the first
two are positive" overlooked --- while the third term is
$\varphi'\cdot(g-f)/\lambda\geq0$, since $\varphi'\geq0$ and $g-f\geq0$ there,
the two agreeing at $\eta_-$ with $g'>f'$. Hence, pointwise,
\[
   h'(\eta)\geq(1-\varphi_\lambda)f'(\eta)+\varphi_\lambda g'(\eta)
   \geq f'(\eta)>0 ,
\]
the middle inequality because $g'>f'$, so $h'$ never vanishes and the collar
has no tangent equal to $u$; and none equal to $-u$ either, a graph over $\eta$ having
$\eta'>0$ throughout. An earlier revision, working from the wrong slope sign,
had both $f'$ and $g'$ negative, needed a Taylor bound on $g-f$ to control the
third term, and chose $\lambda$ small to make the sum negative. With the frame
computed the third term is on the same side as the other two and no choice of
$\lambda$ is required.

\emph{(3) It creates no double point and changes no turning.} Pointwise
$f\leq h\leq g$ on the collar --- the value is a convex combination and $g\geq f$
there, which is the separation the corrected slope chain supplies --- so the
collar lies in the region between the two graphs; that region is inside the
disc $\Delta$ and inside the embedded chart, where the only strands of the
diagram are the old arc --- whose collar stretch is removed --- and the inserted
arc. The collar is embedded by~(1) and meets those two only at its endpoints,
so no double point is created in it, and the model's own double point, at
$|t|=1$, is interior to the inserted arc and untouched. For the turning: the
collar's tangent angle is continuous, stays in $(-\pi/2,0)$ --- the slope
$h'$ being positive and finite, so the tangent angle $\theta$ has
$-\tan\theta=h'>0$ with $\cos\theta>0$ --- and takes at its ends exactly the
angles of the two arcs it joins, so the lift of the tangent
along the modified curve is the concatenation of the old lifts with an
interpolating stretch of the same endpoint values; the total tangent winding is
unchanged by its insertion. An earlier revision stopped at $C^1$
and the document then carried a lemma to recover smoothness. This sentence is what the downstream consumers of this
lemma read, the front construction below taking tangent directions of the
curve, and an intermediate revision of this proof deleted it while rewriting
the tail; it is restored here.

\emph{The collar at $p_+$.} The same three properties are now proved at the
other join, with its own signs printed: here the \emph{incoming} arc is the
inserted one and the \emph{outgoing} arc is the old one, and both slopes are
negative. Let $\eta_+=\eta(p_+)>\eta(p)$. On a stretch
$[\eta_+-\lambda_+,\eta_+]$ chosen above $\eta(p)$ and inside the stretches
on which both arcs are graphs, write $\xi=g(\eta)$ for the inserted arc and
$\xi=f(\eta)$ for the old arc extended backward from $p_+$; the two agree at
$\eta_+$ with common slope $m_+=-\tan\theta_+<0$, the angle
$\theta_+=\theta(s_+)$ lying in $(0,\tfrac\pi2)$. Just left of the join the
inserted arc turns strictly negatively into its endpoint, so its angle
exceeds $\theta_+$ there, while the old arc turns positively, so its angle
lies below $\theta_+$; by $\text{slope}=-\tan\theta$, decreasing in
$\theta$,
\[
   g'(\eta)<m_+<f'(\eta)<0
   \qquad\text{on }[\eta_+-\lambda_+,\eta_+),
\]
the outer inequality $f'<0$ because $\theta>0$ above $\eta(p)$. Integrating
$g'-f'<0$ backward from the common value at $\eta_+$ gives $g>f$ on
$[\eta_+-\lambda_+,\eta_+)$. Set
\[
   h=\bigl(1-\chi_{\lambda_+}\bigr)g+\chi_{\lambda_+}f,
   \qquad
   \chi_{\lambda_+}(\eta)=\varphi\Bigl(\frac{\eta-(\eta_+-\lambda_+)}{\lambda_+}\Bigr),
\]
with the same transition profile $\varphi$: the modified curve runs along the
inserted arc up to $\eta_+-\lambda_+$, along the graph of $h$ across the
collar, and along the old arc from $\eta_+$ on, the two joins $C^\infty$ by
the flat endpoint jets of $\varphi$ exactly as in (1). For the derivative,
\[
   h'=\bigl(1-\chi_{\lambda_+}\bigr)g'+\chi_{\lambda_+}f'
      +\frac{\varphi'\bigl((\eta-\eta_++\lambda_+)/\lambda_+\bigr)}{\lambda_+}
       \,\bigl(f-g\bigr),
\]
and the sign of $f-g$ in the cutoff term is exactly what the negativity of
the two slopes alone does not control: here $f-g\leq0$ by the separation just
integrated, and $\varphi'\geq0$, so the cutoff term is nonpositive, while the
convex part is a convex combination of two negative numbers. Hence,
pointwise,
\[
   h'(\eta)\leq\bigl(1-\chi_{\lambda_+}\bigr)g'(\eta)+\chi_{\lambda_+}f'(\eta)
   \leq f'(\eta)<0 ,
\]
the middle inequality because $g'<f'$, and the bound pointwise for the
mirrored reason: $f'$ varies along the collar and no endpoint bound is
available. So $h'$ never vanishes, and this collar --- a graph over $\eta$
with $\eta'>0$ throughout --- has no tangent equal to $u$ and none equal to
$-u$; its tangent angle stays in $(0,\tfrac\pi2)$, with
$-\tan\theta=h'<0$ and $\cos\theta>0$, and takes at its ends exactly the
angles of the two arcs it joins, so the lift concatenation of (3) applies to
it verbatim and the total tangent winding is unchanged. Pointwise
$f\leq h\leq g$, the value being a convex combination and $g\geq f$ there, so
the collar lies in the region between the two graphs, inside the disc
$\Delta$ and the embedded chart, and it is embedded and meets the two arcs
only at its endpoints: no double point is created, by the argument of (3)
unchanged.

\emph{The four conclusions, separately.} (ii): $c'(t)=(2t,1-3t^2)$ is parallel to
$u_0$ only at $t=0$, where $c'(0)=-u_0$; since $A$ is invertible with
$Au_0=u$, the inserted arc has no tangent $u$ and exactly one tangent $-u$, and
the retained tails have neither, their angles lying in $(-\tfrac\pi2,\tfrac\pi2)$
away from $0$.

(iii): if $c(s)=c(t)$ then $s=\pm t$, and the second coordinate forces
$2t(1-t^2)=0$, so the only pair of distinct parameters is $\{-1,1\}$: exactly
one double point. With the later branch over the earlier one its sign is that of
$\det(c'(1),c'(-1))=-8<0$, and $\det B>0$ preserves it.

(iv) for $\rot$: Let $R$ be the rotation carrying the positive orthonormal basis
$(v_0,u_0)$ to $(v,u)$ and put $C=R^{-1}A$. Relative to $(v_0,u_0)$,
$C$ has matrix
$\left(\begin{smallmatrix}x&0\\y&1\end{smallmatrix}\right)$ with $x>0$, so
$C_t=\left(\begin{smallmatrix}1-t+tx&0\\ty&1\end{smallmatrix}\right)$
joins $I$ to $C$ inside $\mathrm{GL}^+(2,\mathbb R)$. If $R_t$ rotates
through $t$ times any fixed angle representing $R$, first $C_t$ and then
$R_tC$ give a path $A_t$ from $I$ to $A=RC$ in that group. Apply
$z\mapsto A_tz/|A_tz|$ to the closed direction loop obtained by following
the tangent path of $c$ and the tangent path of $b$ backwards.
Lemma~\ref{lem:turnlift}(iii) therefore preserves their compatible endpoint
lift-increment difference. The tangent of $b$ follows the short positive arc between the
endpoint rays; the tangent of $c$ follows the complementary clockwise arc, since
$\det(c',c'')=-2(1+3t^2)<0$ and it passes through $-u_0$ at $t=0$; so their
relative winding is $-1$. The old central arc also follows the short positive
arc, its lifted sweep $\theta(s_+)-\theta(s_-)$ lying in $(0,\pi)$. Hence Lemma~\ref{lem:turnlift}(iii) says that the replacement changes the
closed curve's rotation by exactly $-1$. For
$w$: the removed arc was embedded and carried no crossing, so the writhe changes
by the one new negative crossing, that is by $-1$.

\emph{The polynomial conclusion, without a knot-category detour.}
The construction just completed also proves that $F'$ remains in the lemma's
input class.  It replaces one oriented arc by one $C^\infty$ regular oriented
arc, leaves the component count unchanged, and creates exactly one transverse
double point in a disc meeting no other strand.  Thus its double points remain
finite and transverse and no triple point is created.

Let $F^{\circ}$ be the oriented diagram obtained from $F'$ by deleting the
unique monogon in $\Delta$ by a Reidemeister-I move.  The move is available:
the preceding separation argument shows that the monogon is disjoint from
all other strands, and the crossing calculation shows it is the unique new
double point.  It replaces that monogon by one crossing-free regular arc in
$\Delta$, so $F^{\circ}$ is again a one-component immersed-circle diagram with
finite transverse double points and no triple points.  Reidemeister-I
invariance in Axiom~\ref{ax:homfly} gives $P_{F'}=P_{F^{\circ}}$.

Outside $\Delta$, the diagrams $F^{\circ}$ and $F$ coincide.  Inside
$\Delta$, each has one crossing-free oriented arc joining the same two
boundary points in the same traversal direction.  Match every marked
preimage outside that local traversal interval with the identical marked
preimage of the other diagram.  The local intervals contain no marked
preimage, so this bijection preserves cyclic order, crossing pairs,
over/under positions and signs: it is a record isomorphism.  Both underlying
curves are one-component immersed circles with transverse double points and
no triple points, so Axiom~\ref{ax:gausscode} says that the diagrams present
the same oriented link.  Axiom~\ref{ax:homfly} therefore gives
$P_{F^{\circ}}=P_F$, hence $P_{F'}=P_F$.  Equality outside $\Delta$ proves the
remaining clause of~(i).
\end{proof}

\begin{remark}[authorship of the construction above, and what its referee
found]\label{rem:curlauthor}
The proof of Lemma~\ref{lem:curl} from the ordered-ray lemma through the
geometric insertion is \textbf{not this author's}.  Its earlier form was
refuted by the peer: it swapped names in one ray triple, left the affine
shrinking unquantified, and did not exclude additional intersections.  The
printed repair is the peer's: the positive-turn chart, monotone equal-level
cuts, explicit $H,x,y$ affine fit, the bound
$1-(qy)^2=4abcd/(bc+ad)^2$, the transverse-excess separation
$\frac{\ell x}{12}(4-t^2)>0$, and the separate tangent, crossing and
rotation conclusions.

The HOMFLY conclusion is narrower than the former knot-isotopy conclusion.
It uses only Reidemeister-I invariance from Axiom~\ref{ax:homfly} and the
record comparison from Axiom~\ref{ax:gausscode}; no standard realization,
arc-isotopy theorem or category-change theorem is involved.

Theorem~\ref{thm:carrierfloor}(C)'s turn hypothesis was tightened for the same
reason: as printed it read ``every turn positive, or exactly one negative'',
which does not exclude a zero turn, while Lemma~\ref{lem:rounding} needs every
turn nonzero. The hypothesis now requires nonzero turns explicitly and says
``exactly one is negative and every other is positive''. Every consumer
supplies them: a carrier's corners are
guarded by (G1) at a vertex and (G5) at a smoothing site
(Definition~\ref{def:wind}).
\end{remark}

\subsection{Global orientation reversal}



\subsection{The floor}

\begin{theorem}[global reversal, rounding record, vertical tangencies and the carrier floor]\label{thm:carrierfloor}
\begin{enumerate}
\item[(R)] Here $\rot$ is as in Definition~\ref{def:rot}.
Let $K$ be an oriented knot and $-K$ the same knot with its orientation
reversed. Then $P_{-K}=P_{K}$. Moreover a diagram of $-K$ obtained by reversing
the orientation of a diagram of $K$ has the same crossing signs, the same
writhe, and rotation number of the underlying plane curve negated.
\item[(A)] Let $L$ and $D$ satisfy the hypotheses of Lemma~\ref{lem:rounding}. For every
$\varepsilon\in(0,\varepsilon_0(L))$, the construction in the proof of that
lemma returns \emph{one} curve and \emph{one} diagram at those data: the
junction inserted at the corner $q_i$ is determined by $\varepsilon$, by the
two incident unit directions and by the transition profile, which is fixed once
and for all in that proof and is not chosen per corner or per curve; the arc
length $\ell$ is then determined by the endpoint condition, and the rest of the
curve is $L$ itself. Write
\[
   \mathrm{Round}(L,D,\varepsilon)=(L_\varepsilon,D_\varepsilon)
\]
for that pair, the \emph{rounding record} of $(L,D)$ at $\varepsilon$, and call
$L_\varepsilon$ the rounded curve and $D_\varepsilon$ the rounded diagram. The
record exists to be quantified over: an earlier revision, written when the
profile was a per-use choice, had every consumer say "every curve the lemma
produces", which is a quantifier over an unnamed family and reads differently
in each consumer.
\item[(B)] Here $\rot$ is as in Definition~\ref{def:rot}.
Let $L$ be a diagrammatic closed polygon (Definition~\ref{def:diagrammatic})
whose principal turns all exist and are nonzero, carrying a diagram $D$, and
assume, after reversing the orientation of $L$ if necessary, that either every
principal turn is positive, or exactly one is negative. Put $R=|\rot(L)|$. Then
there are a direction $u\in S^1$ and an $\varepsilon_1>0$ such that \emph{for
every} $\varepsilon\in(0,\varepsilon_1)$ the rounded curve $L_\varepsilon$ of
the record $\mathrm{Round}(L,D,\varepsilon)$
from clause~(A) has exactly $R$ points at which its unit
tangent equals $u$ and exactly $R$ at which it equals $-u$; at each of them the
tangent crosses that direction in the positive sense. The hypotheses are stated
here rather than imported by "as in Lemma~\ref{lem:rounding}", and the
clearance is produced here rather than borrowed from that lemma's existential
one. The object is named rather than quantified over: the record is one curve
and one diagram at each $\varepsilon$, the profile being fixed once and for all,
and an earlier revision said "every curve the lemma produces" because at the
time the profile was a per-use choice.
\item[(C)] Here $\rot$ is as in Definition~\ref{def:rot}.
Let $D$ be an oriented knot diagram all of whose crossings are positive, with
writhe $w$, whose underlying plane curve is a closed polygon $L$ with all
principal turns existing, \emph{nonzero}, and of magnitude below $\pi$, with
finitely many double points, all transversal, with no triple points, none of
them a corner of $L$, and no corner of $L$ lying on a non-incident edge; and let
$R$ be the absolute value of its Whitney rotation number.  The
finiteness/transversality and the two corner conditions are what
Lemma~\ref{lem:rounding} asks of its input.  The no-triple condition repeats the
shared-image clause of Definition~\ref{def:diagrammatic}; it is preserved by
rounding and is needed below when diagram records are compared.  These
conditions are stated here so that no generic parent polygon is assumed. Assume that, after reversing the
orientation if necessary --- which by clause~(R) changes neither
$P_D$, nor the crossings and their signs, nor $w$, nor $R$ --- either every
principal turn is positive, or exactly one is negative and every other is
positive. Then
\begin{equation}
\min\deg_a P_D(a,z)\;\geq\;1-w-R,
\label{eq:floor}
\end{equation}
and the same bound holds for $f_D(a)=[z^0]P_D(a,z)$ whenever $f_D\neq0$.
\item[(D)] Here $\rot$ is as in Definition~\ref{def:rot}.
Let $P$ be generic, let $S\in\Ind(G_P)$ and let $L$ be a carrier of $S$ carrying
no residual piece, so that $P_{S,L}=1$ and $w_{S,L}=0$ by
Definition~\ref{def:X1}'s empty conventions --- an earlier revision wrote those
two symbols with $S$ never introduced --- and suppose that all its principal turns are nonzero and that,
after reversing the orientation if necessary, either every turn is positive or
exactly one is negative. Then $R(L)\geq1$ and
$\min\deg_a P_{S,L}=0\geq 1-w_{S,L}-R(L)$.
\end{enumerate}
\end{theorem}
\begin{proof}
\emph{Clause (R).}
Reversing the orientation of all components of a link sends a skein triple
$(L_+,L_-,L_0)$ to a skein triple: the sign of a crossing is
$\sgn\det(u_{\mathrm{over}},u_{\mathrm{under}})$ and reversing both strands
replaces both vectors by their negatives, leaving the determinant unchanged, and
the oriented smoothing of the reversed diagram is the reverse of the oriented
smoothing. So $L\mapsto P_{-L}$ satisfies the defining skein of
Axiom~\ref{ax:homfly}, and it sends the unknot to $1$; by the uniqueness clause
it equals $L\mapsto P_L$. The diagram statements have the same scope. Every crossing sign is a
determinant of two reversed vectors, so the writhe is unchanged. For a
polygonal underlying curve every principal turn changes sign and
Lemma~\ref{lem:turnlift}(ii) negates rotation. For a $C^1$ regular curve,
if $\theta$ lifts its normalised tangent, then $\theta(1-s)+\pi$ lifts the
reversed tangent with the negative endpoint increment, as in
Lemma~\ref{lem:turnlift}(i).

\emph{Clause (A).} The transition profile is fixed once and for all in the proof of
Lemma~\ref{lem:rounding}.  At each corner, $\varepsilon$ and the two incident
unit directions determine that profile's junction, and the endpoint condition
determines its length $\ell$; outside the rounding discs the curve is $L$, and
the diagram data are inherited.  Thus the named construction returns one
specified pair at every allowed $\varepsilon$.  This asserts uniqueness of the
fixed construction's record, not uniqueness among all possible smoothings.

\emph{Clause (B).}
Reversing the orientation of $L$ negates every principal turn and the signed rotation and
preserves every self-crossing sign, so perform the allowed normalization once. In the
uniform case $\rot(L)=R\geq1$ by Lemma~\ref{lem:uniformrot}(i); in the one-dissent case the same holds by Lemma~\ref{lem:uniformrot}(ii).

Take $\varepsilon_1=\varepsilon_0(L)$, the clearance of Lemma~\ref{lem:rounding}.
Choose $u\in S^1$ so that neither $u$ nor $-u$ is an edge direction and, in the
one-dissent case, neither lies in the negative turn's closed swept arc $A$, of
length $\alpha:=|\tau_-|<\pi$. Such $u$ exists: in the one-dissent case $A\cup(-A)$
has length at most $2\alpha<2\pi$, while in the uniform case there is no such arc; the additional
forbidden set of edge directions and their antipodes is finite. This choice depends only on $L$. Fix an argument $\phi$ of $u$.

Fix $\varepsilon\in(0,\varepsilon_1)$, write
$(L_\varepsilon,D_\varepsilon)=\mathrm{Round}(L,D,\varepsilon)$, and let $T$ be
its unit-tangent map. Parameterize the rounded traversal by $s\in[0,1]$, with
$s=0$ in the interior of a straight part; because $u,-u$ are not edge directions, its tangent
is congruent to neither $\phi$ nor $\phi+\pi$. Concatenating straight-part arguments
with the junction lifts supplied by Lemma~\ref{lem:rounding}(b) gives a continuous
$\theta\colon[0,1]\to\mathbb R$, starting at an argument $\theta(0)$ of $T(0)$.
Since $T(1)=T(0)$, write $\theta(1)=\theta(0)+2\pi m$ for an integer $m$.

Choose a direction $r$ not parallel to any edge, with argument $\gamma$. The
change across junction $i$ of $\lfloor(\theta-\gamma)/(2\pi)\rfloor$ is $+1$,
$-1$ or $0$ in exactly the three determinant cases defining $\epsilon_i$ in
Definition~\ref{def:rot}. Indeed the same clause makes the lifted change
$\tau_i$, and $|\tau_i|<\pi$ permits at most one crossed level, in its
signed direction; there is no change on a straight part. Telescoping and
$\lfloor q+m\rfloor-\lfloor q\rfloor=m$ give
$m=\sum_i\epsilon_i=\rot(L)=R$. Thus
\[
 \theta(1)=\theta(0)+2\pi R.                                  \tag{*}
\]

Fix $\beta=\phi$ or $\phi+\pi$. Because $u,-u$ are not edge directions, no
level $\beta+2\pi k$ is met on a straight part or at a junction endpoint; by
the choice of $u$, none is met on the decreasing junction. Hence floor jumps
occur only in increasing junctions and never elsewhere. Lemma~\ref{lem:rounding}(b)
makes each occurrence unique with positive angular derivative. Put
\[
 x=\frac{\theta(0)-\phi}{2\pi},\qquad
 y=\frac{\theta(0)-(\phi+\pi)}{2\pi}.
\]
For $u$, (*) and the integer floor identity give, separately,
\[
 \#T^{-1}(u)=\lfloor x+R\rfloor-\lfloor x\rfloor=R.
\]
For $-u$, the same argument at the other levels gives
\[
 \#T^{-1}(-u)=\lfloor y+R\rfloor-\lfloor y\rfloor=R.
\]
There are finitely many junctions, so the preimages are finite; their positive derivatives prove the assertion for arbitrary $\varepsilon$.

\emph{Clause (C).} Fix $\varepsilon\in(0,\varepsilon_1)$ with $\varepsilon_1$ the clearance of
clause~(B) and pass to the rounding record
$(L_\varepsilon,D_\varepsilon)=\mathrm{Round}(L,D,\varepsilon)$ of
clause~(A) --- one curve and one diagram, not a member
of an unnamed family. By Lemma~\ref{lem:rounding}(c),(d) the diagram, its
crossings, its writhe and its rotation number are unchanged, and by
clause~(B) there is a direction $u$ met by the tangent of
$L_\varepsilon$ exactly $R$ times, each time positively, and likewise for
$-u$. Rotate coordinates so
that $u$ is the downward vertical. The rounded diagram then has exactly $R$
downward vertical tangencies and $R$ upward ones.

Switch every crossing. The result is a diagram $\overline D$ of the mirror knot,
with the same underlying plane curve --- hence the same tangencies --- and with
writhe $-w$; every crossing of $\overline D$ is negative. Apply
Lemma~\ref{lem:curl} at each of the $R$ downward vertical tangencies; they are
isolated, lie in the interiors of positively turning rounding arcs
(clause~(B)), and the rounding discs are pairwise disjoint and
meet no crossing, by the choice of $\varepsilon$ in the proof of
Lemma~\ref{lem:rounding}, so the $R$ discs are
disjoint and each replacement leaves the others intact. Call the result $T$. By
Lemma~\ref{lem:curl}, $P_T=P_{\overline D}$, $T$ has no downward vertical
tangency, and it has writhe
\[
w(T)=-w-R .
\]
Every crossing of $T$ is negative: the old ones by the switch, the $R$ new ones by
Lemma~\ref{lem:curl}(iii).

We now construct the transverse lift rather than import a
front criterion.  Write the smooth regular parametrisation of $T$ as
$s\mapsto(x(s),z(s))$, put
\[
 v(s)=\sqrt{x'(s)^2+z'(s)^2},\qquad
 y_0(s)=-\frac{x'(s)}{v(s)+z'(s)},
\]
and use the convention that, in the $xz$ projection, the branch with smaller
$y$-coordinate is drawn over.  Since $T$ has no downward vertical tangency,
$v+z'>0$ everywhere: equality would force $x'=0$ and $z'<0$.  Thus $y_0$ is
smooth and periodic, and direct algebra, with no division by $x'$, gives
\[
 z'-y_0x'=z'+\frac{x'^2}{v+z'}
 =z'+\frac{v^2-z'^2}{v+z'}=v>0.                 \tag{*}
\]

It remains only to impose the crossing order.  At a crossing branch there is
no vertical tangency, so put $m=z'/x'$.  The strict inequality in~(*) permits
$y<m$ on a rightward branch and $y>m$ on a leftward branch.  For every crossing
choose admissible constants $c_{\rm O}<c_{\rm U}$ at its over- and underpassing
preimages.  Such a choice is automatic if the overpass points right: after
choosing any admissible $c_{\rm U}$, take
$c_{\rm O}<\min\{m_{\rm O},c_{\rm U}\}$.  If both point left, choose any
$c_{\rm O}>m_{\rm O}$ and then
$c_{\rm U}>\max\{m_{\rm U},c_{\rm O}\}$.  The only remaining case has
$x'_{\rm O}<0<x'_{\rm U}$.  Every crossing of $T$ is negative, and here
\[
 0>\det(t_{\rm O},t_{\rm U})
   =x'_{\rm O}x'_{\rm U}(m_{\rm U}-m_{\rm O}).
\]
Since $x'_{\rm O}x'_{\rm U}<0$, this says $m_{\rm O}<m_{\rm U}$; choose
$m_{\rm O}<c_{\rm O}<c_{\rm U}<m_{\rm U}$.  These are exactly the two allowed
half-lines and the required over/under order.

There are finitely many crossing preimages.  Around each one $s_j$, choose
pairwise disjoint cyclic parameter intervals $(\alpha_j,\beta_j)$ so small that
$z'-c_jx'>0$ throughout.  Using the explicit flat transition profile
$\varphi$ displayed in the proof of Lemma~\ref{lem:rounding}, define a smooth
bump on the circle by
\[
 \psi_j(s)=
 \begin{cases}
 \varphi\!\left((s-\alpha_j)/(s_j-\alpha_j)\right),&
       \alpha_j\le s\le s_j,\\
 \varphi\!\left((\beta_j-s)/(\beta_j-s_j)\right),&
       s_j\le s\le\beta_j,\\
 0,&\text{otherwise}.
 \end{cases}
\]
All endpoint derivatives are zero, so this is smooth, equals one at $s_j$ and
has support in that interval.  Set
\[
 y=y_0+\sum_j\psi_j(c_j-y_0).
\]
The supports are disjoint.  On the $j$-th one,
\[
 z'-yx'=(1-\psi_j)(z'-y_0x')+\psi_j(z'-c_jx')>0,
\]
and outside them this is~(*).  Hence $y$ is smooth and periodic and the lift
$s\mapsto(x(s),y(s),z(s))$ is positively transverse to $dz-y\,dx$.

Finally the lift is regular because its $xz$ projection is immersed.  If two
lifted points agreed, their projected points would agree; the only distinct
parameters with that property are the two preimages of a listed double point,
and there $c_{\rm O}<c_{\rm U}$.  Thus the lift is injective, hence, by
compactness of the parameter circle, an embedded transverse knot.  Its front
and its smaller-$y$ over/under assignments are exactly $T$.

By Axiom~\ref{ax:etnyre}, $\operatorname{sl}(T)=w(T)=-w-R$. Since
$P_T=P_{\overline D}$, Axiom~\ref{ax:slbound} applied to $T$ gives
\[
\max\deg_a P_{\overline D}\;\leq\;-\operatorname{sl}(T)-1\;=\;w+R-1 .
\]
For every oriented-link isotopy class $K$, put
$\mathcal M(K)=P_{\overline K}(a,z)$. Mirroring exchanges $L_+$ with $L_-$
and fixes $L_0$, so the defining skein of Axiom~\ref{ax:homfly} on the
mirrored triple, multiplied by $-1$, is
\[
 a^{-1}\mathcal M(L_+)-a\mathcal M(L_-)=-z\,\mathcal M(L_0),
 \qquad \mathcal M(\bigcirc)=1.
\]
The substitution $(a,z)\mapsto(a^{-1},-z)$ is involutive, so applying it to
the uniqueness clause of Axiom~\ref{ax:homfly} gives uniqueness for these two
conditions. The assignment $K\mapsto P_K(a^{-1},-z)$ satisfies them; hence
$P_{\overline K}(a,z)=P_K(a^{-1},-z)$. As a Laurent-ring automorphism, the
substitution negates every $a$-exponent and cannot cancel a nonzero extreme
coefficient. Hence
$\max\deg_a P_{\overline D}=-\min\deg_a P_D$, and the displayed bound becomes
$\min\deg_a P_D=-\max\deg_a P_{\overline D}\geq1-w-R$, which is
\eqref{eq:floor}. Finally, the support of $f_D$ is contained in the support of
$P_D$, so the same bound holds for it when it is nonzero.

\emph{Clause (D).} In the uniform case $|\rot(L)|\geq1$ by Lemma~\ref{lem:uniformrot}(i); in the
one-dissent case $\rot(L)\geq1$ by Lemma~\ref{lem:uniformrot}(ii). Either way
$R(L)\geq1$, so $1-0-R(L)\leq0=\min\deg_a1$. In particular the hypothesis is met by every carrier of a live support
(Definition~\ref{def:wind}), and it is met by the marked halves of the
one-dissent geometry produced in the vertex-on-edge section, where that is
proved at the point of use.
\end{proof}

\begin{remark}[what the floor does and does not assert]\label{rem:floorscope}
Theorem~\ref{thm:carrierfloor}(C) is a lower bound on a minimal degree; it asserts
no sharpness. For a live carrier the slot of Definition~\ref{def:X1} is exactly
$1-w_{S,L}-R(L)$, so exactly one of two things holds: the slot is the bottom
supported degree of the row, or the slot lies strictly below the support and its
coefficient is zero. Both are consistent with \eqref{eq:floor} and the proofs of
Section~\ref{sec:s7} use only the second reading, namely that nothing lies below
the slot.
\end{remark}
